Generated explanations can remain plausible after their evidence changes.
We study a typed temporal property graph that records observations, proposed
propositions, declared prerequisites and support at each knowledge time.
A query-specific semantic core retains numerical witnesses and version history.
A shared network view lets readers inspect original citations, corrected versions
and alternative support routes in the paper and dashboard. Two conditional results
bound which claims may change and characterize when direct maintenance misses
required updates. We measure these misses through transitive exposure, the fraction
of necessary changes excluded by direct scheduling. In synthetic, WESAD and
PPG-DaLiA experiments, only \SemParentClaims{} of \SemDisplayedClaims{} displayed
claims declare parents. Every eligible waveform revision has zero exposure.
Independent scoring separates source grounding from task completeness.
Checked displays are fully grounded, but subject-averaged required recall is
\SemPpgDisplayR--\SemSyntheticDisplayR\%. Across \SemCoreCases{} controlled
structural cases, exact and GPU-generated programs expose necessary intermediate
dependencies and surviving alternatives. Exposure predicts direct-maintenance
misses in all \SemPredictionGroups{} eligible replays. Full maintenance agrees
across Neo4j and an indexed relational control. Generated answers reach
evidence-to-answer depth one or two despite requests up to four.
These results explain when explicit semantic structure enables selective revision.
Storage size, path length and contract acceptance alone do not establish
explanation fidelity or maintenance benefit.
